% ============================================================================
% WIP — Conclusiones. NO incluidas en el informe todavía (enfoque actual:
% fases 1--3 de CRISP-DM). Reactivar en main.tex cuando se cierre la evaluación.
% ============================================================================
\section{Conclusiones}
\label{sec:conclusiones}

\begin{itemize}
  \item Se construyó un pipeline reproducible de recolección y normalización
        con resultados tipados y trazabilidad por corrida
        (\texttt{data/manifest.json}), cubriendo las dos fuentes del proyecto: la API
        dolarbluebolivia.click y la Gaceta Oficial \cite{dolarblue2026api}.
  \item Se dejó preparada (fases 1--3 de CRISP-DM) la base para predecir el
        nivel del paralelo y la brecha a horizontes de 1--7 días, con
        partición temporal y línea base naïve \cite{hyndman2021forecasting}.
  \item Este informe es, a su vez, parte del entregable: se regenera
        íntegramente (figuras, tablas, diagramas y PDF) con un único comando
        (comando en el Apéndice) desde los últimos CSV disponibles.
\end{itemize}

\subsection{Trabajo futuro}

\begin{itemize}
  \item Modelar y evaluar la previsión (Objetivo 1) y la relación entre la
        Gaceta y la brecha (Objetivo 2) con validación walk-forward
        \cite{box2015timeseries}.
  \item Respetar la cuota semanal de la API ($\sim$30 llamadas por corrida):
        ingestas programadas, no exploratorias.
\end{itemize}
